\chapter{Cel i zakres pracy}

\section{Cel pracy}
Głównym celem projektu było zaprojektowanie i zaimplementowanie aplikacji mobilnej, która w sposób stabilny i wydajny pobiera dane pomiarowe, umożliwia ich przeglądanie w formie historycznego dziennika oraz oferuje zaawansowane opcje personalizacji parametrów sieciowych i wizualnych.
\section{Cele szczegółowe}
Realizacja projektu wymagała skupienia się na kilku kluczowych aspektach technicznych, które pozwoliły stworzyć kompletne narzędzie diagnostyczne. Do najważniejszych celów szczegółowych zaliczono następujące punkty:
\begin{itemize}
    \item opracowanie responsywnego interfejsu użytkownika zgodnego z najnowszymi wytycznymi Material Design.,
    \item implementacja asynchronicznego modułu komunikacji sieciowej opartego na sprawdzonej bibliotece Retrofit,
    \item zaimplementowanie mechanizmu SharedPreferences do trwałego przechowywania konfiguracji użytkownika w pamięci urządzenia,
    \item stworzenie systemu dynamicznej zmiany motywu graficznego pozwalającego na wybór między trybem jasnym a ciemnym,
    \item opracowanie zaawansowanego systemu obsługi błędów komunikacji, który informuje użytkownika o problemach w sposób niezakłócający pracy programu.
\end{itemize}

\section{Zakres pracy}
Zakres pracy został ograniczony do warstwy klienckiej, co pozwoliło na głęboką optymalizację kodu pod kątem wydajności na urządzeniach o różnej specyfikacji sprzętowej. Poniżej przedstawiono kluczowe obszary, które zostały objęte procesem deweloperskim:
\begin{itemize}
    \item projektowanie i implementacja aktywności głównej oraz ekranu ustawień systemowych,
    \item tworzenie adapterów do obsługi dynamicznych list danych w komponencie RecyclerView,
    \item konfiguracja usług sieciowych i mapowanie obiektów JSON na klasy języka Java,
    \item zarządzanie zasobami graficznymi, w tym tworzenie inteligentnych ikon wektorowych zmieniających kolor,
    \item przeprowadzenie testów funkcjonalnych sprawdzających poprawność zapisu i odczytu preferencji użytkownika.
\end{itemize}